% ============================================================
% CAPITOLO N - DAL BASSO NUMERATO AL BASSO SENZA NUMERI
% (Capitolo ponte - NUOVO)
% ============================================================
\chapter{Dal basso numerato al basso senza numeri}
\label{chap:dal_numerato_al_senza_numeri}

Fino a questo punto del manuale, ogni esercizio di armonizzazione ci è stato offerto con un prezioso aiuto: la \textbf{numerica}, ovvero le cifre scritte sotto le note del basso che ci dicono quale accordo usare in ogni momento. La numerica è la bussola del basso figurato: indica lo stato degli accordi, segnala le dissonanze, guida la condotta delle voci.

Arrivare al basso senza numeri non significa buttare via quella bussola. Significa \textbf{averla interiorizzata} al punto da non averne più bisogno: l'obiettivo di questo capitolo è costruire la capacità di dedurre la numerica da sé, prima ancora di scrivere una nota.

\section{Che cosa cambia senza la numerica}
\label{sec:cosa_cambia}

Nel basso \textbf{figurato} (con numeri), il percorso operativo è:
\begin{center}
	Numerica $\rightarrow$ Accordo $\rightarrow$ Condotta delle voci
\end{center}

Nel basso \textbf{senza numeri}, il percorso diventa:
\begin{center}
	Analisi del basso $\rightarrow$ Numerica implicita $\rightarrow$ Accordo $\rightarrow$ Condotta delle voci
\end{center}

Sembra una piccola differenza. In realtà, il primo passo --- ``analisi del basso'' --- richiede la mobilitazione di \textit{tutte} le conoscenze acquisite: tonalità, funzioni armoniche, cadenze, progressioni, modulazioni, ritmo armonico. È un'operazione di \textbf{sintesi} che presuppone la padronanza di tutti i nuclei analitici precedenti.

\begin{osservazione}[Il basso senza numeri è un'operazione di scelta responsabile]
	\label{oss:scelta_responsabile}
	Nel basso numerato, lo studente applica regole note a una situazione determinata. Nel basso senza numeri, lo studente \textbf{sceglie} quali regole applicare e in quale contesto. Questa differenza è profonda: non si tratta di sapere di più, ma di usare quello che si sa in modo \textit{autonomo} e \textit{argomentato}.
\end{osservazione}

\section{La lettura del basso come testo musicale}
\label{sec:lettura_basso}

Il basso senza numeri è una \textbf{linea melodica} che contiene, implicite, informazioni armoniche. Saperlo leggere significa saper estrarre quelle informazioni prima ancora di sedersi al pianoforte o al foglio di carta.

\subsection{Informazioni immediate: tonalità e modo}

La prima operazione è sempre l'identificazione della \textbf{tonalità e del modo}.

\begin{regola}[Identificazione della tonalità nel basso senza numeri]
	\label{reg:identificazione_tonalita}
	Per determinare la tonalità e il modo di un basso senza numeri, si seguono questi passi:
	\begin{enumerate}
		\item \textbf{Armatura di chiave:} è il punto di partenza. Stabilisce il pool di note ``naturali'' del brano.
		\item \textbf{Nota iniziale e finale:} tipicamente la nota iniziale è la tonica o la dominante; la nota finale è quasi sempre la tonica. Questo conferma (o smentisce) l'ipotesi suggerita dall'armatura.
		\item \textbf{Alterazioni cromatiche:} la presenza di una sensibile alzata (VII grado) conferma il \textit{modo minore armonico}. Alterazioni che non appartengono alla scala di riferimento segnalano tonicizzazioni o modulazioni (cfr. \ref{chap:tonicizzazione} e \ref{chap:modulazione}).
		\item \textbf{Salto di quarta o quinta iniziale:} se il basso inizia con un salto I-V o V-I, la tonalità è subito chiara.
	\end{enumerate}
\end{regola}

\begin{esempio}[Identificazione della tonalità da un basso senza numeri]
	\label{ex:identificazione_tonalita}
	Basso da analizzare in Do minore:
	\musica{
		\new Staff {
			\clef bass \key do \minor \time 4/4
			\relative do, {
				do4 re mib fa | sol2 sold4 la | sib2 la4 sol | do1
			}
		}
	}
	Osservazioni:
	\begin{itemize}
		\item Armatura: 3 bemolli $\rightarrow$ Do minore (o \Mibemolle  maggiore).
		\item Nota iniziale e finale: Do $\rightarrow$ conferma Do minore.
		\item \Soldiesis : sensibile alzata $\rightarrow$ modo minore armonico confermato.
		\item Il \Ladiesis  in realtà è un La$\natural$ (la sensibile alzata è il Sol\diesis{}) --- in questo basso il La è la sopradominante naturale.
	\end{itemize}
\end{esempio}

\subsection{Informazioni strutturali: cadenze e frasi}

Una volta stabilita la tonalità, il secondo passo è la \textbf{segmentazione fraseologica}: trovare dove si trovano le cadenze e quindi dove sono i confini tra le frasi.

\begin{regola}[Riconoscimento delle cadenze nel basso senza numeri]
	\label{reg:riconoscimento_cadenze_basso}
	Le cadenze si riconoscono cercando le seguenti \textbf{formule armoniche tipiche} alla fine delle frasi:
	\begin{itemize}
		\item \textbf{Cadenza perfetta} (V-I): il basso sale di quarta (o scende di quinta) da V a I, entrambi in stato fondamentale.
		\item \textbf{Cadenza sospesa} (--V): il basso termina sulla nota della dominante, spesso preceduta da IV o II.
		\item \textbf{Cadenza plagale} (IV-I): il basso scende di quarta da IV a I.
		\item \textbf{Cadenza composta} (IV-V-I o II-V-I): tre note del basso che formano la formula conclusiva più completa.
		\item \textbf{Cadenza d'inganno} (V-VI): il basso sale di grado congiunto da V a VI dopo aver fatto aspettare la tonica.
	\end{itemize}
	Le cadenze si trovano quasi sempre ai bordi delle frasi di 4 misure (struttura antecedente-conseguente, cfr. \ref{chap:struttura_frase}).
\end{regola}

\subsection{Informazioni dinamiche: il ritmo armonico}

Il \textbf{ritmo armonico} è la frequenza con cui gli accordi cambiano all'interno di una battuta. È un'informazione cruciale che spesso il basso ci fornisce implicitamente.

\begin{definizione}[Ritmo armonico]
	\label{def:ritmo_armonico}
	Il \textbf{ritmo armonico} è la durata media di ogni accordo all'interno del brano. Un ritmo armonico \textit{lento} significa accordi lunghi (semibrevi o minime); un ritmo armonico \textit{veloce} significa accordi corti (semiminime o crome).
\end{definizione}

\begin{regola}[Ritmo armonico e numero di accordi per battuta]
	\label{reg:ritmo_armonico_battuta}
	In un basso senza numeri in stile corale:
	\begin{itemize}
		\item In \textbf{4/4} e \textbf{C} (alla breve): tipicamente uno o due accordi per battuta (uno per minima).
		\item In \textbf{3/4}: tipicamente uno o tre accordi per battuta.
		\item In \textbf{2/2} o \textbf{3/2}: tipicamente un accordo per minima.
		\item Ogni nota del basso che vale più di un tempo \textit{può} essere armonizzata con uno o più accordi diversi --- ma il cambio armonico deve essere argomentato.
		\item Una nota di basso che vale un solo movimento (semiminima in 4/4) di solito porta un accordo solo.
	\end{itemize}
	\textbf{Regola pratica di default:} in un esercizio scolastico, si assegna un accordo per ogni \textit{accento metrico principale} della misura, a meno che il contesto non indichi chiaramente il contrario.
\end{regola}

\section{Il processo deduttivo: da nota a funzione}
\label{sec:processo_deduttivo}

Il cuore del lavoro sul basso senza numeri è il \textbf{ragionamento funzionale}: partendo dalla nota del basso, dedurre quale funzione armonica essa esercita nel contesto in cui si trova.

\subsection{La nota del basso non è una funzione isolata}

Questo è il principio fondamentale, evidenziato anche dalle fonti metodologiche di riferimento: \textbf{una nota del basso, presa isolatamente, non esprime alcuna funzione tonale}. Il suo significato diventa chiaro solo nel confronto con le note vicine, nel contesto delle cadenze e delle progressioni riconoscibili.

\begin{esempio}[Ambiguità funzionale di un Sol nel basso]
	\label{ex:ambiguita_sol}
	La nota Sol al basso può essere:
	\begin{itemize}
		\item \textbf{V grado in Do maggiore} (dominante) $\rightarrow$ accordo di Sol maggiore (V)
		\item \textbf{I grado in Sol maggiore} (tonica) $\rightarrow$ accordo di Sol maggiore (I)
		\item \textbf{IV grado in Re maggiore} $\rightarrow$ accordo di Sol maggiore (IV) o Sol settima
		\item \textbf{VII grado in La minore} (sensibile alzata) $\rightarrow$ accordo di \Soldiesis  diminuito (VII°)
		\item \textbf{Nota di passaggio} tra La e Fa $\rightarrow$ nessun cambio armonico
	\end{itemize}
	Solo il contesto --- le note precedenti e seguenti, le alterazioni, la posizione metrica --- ci dice quale delle interpretazioni è corretta.
\end{esempio}

\subsection{Gerarchia delle informazioni contestuali}

Quando si analizza una nota del basso, le informazioni utili vanno cercate in questo ordine di priorità:

\begin{regola}[Gerarchia delle informazioni per la deduzione armonica]
	\label{reg:gerarchia_informazioni}
	\begin{enumerate}
		\item \textbf{Posizione metrica:} le note sui tempi forti portano quasi sempre accordi in stato fondamentale (cfr. \ref{reg:stato_basso}). Le note sui tempi deboli possono portare rivolti o note di passaggio.
		\item \textbf{Presenza di alterazioni:} un'alterazione cromatica segnala quasi sempre una funzione speciale: sensibile del modo minore, dominante secondaria, o cambio di tonalità.
		\item \textbf{Movimento rispetto alla nota precedente:} grado congiunto suggerisce Regola dell'Ottava o progressione; salto di quarta/quinta suggerisce cambio di pilastro armonico; salto di ottava può indicare cambio di rivolto.
		\item \textbf{Contesto cadenzale:} se la nota si trova in prossimità di un punto di cadenza già individuato, la sua funzione è determinata dalla formula cadenzale.
		\item \textbf{Progressioni riconoscibili:} se le note del basso formano un pattern di progressione (cfr. \ref{chap:progressioni_armoniche}), la funzione è determinata dalla struttura sequenziale.
	\end{enumerate}
\end{regola}

\section{Applicazione della Regola dell'Ottava al basso senza numeri}
\label{sec:regola_ottava_basso_senza_numeri}

La \textbf{Regola dell'Ottava} (cfr. Capitolo \ref{chap:funzioni_armoniche} e \ref{chap:regola_ottava_completa}) è lo strumento principale per armonizzare i movimenti scalari del basso. Nel basso senza numeri, si applica la stessa logica, ma la scelta del livello di complessità (versione preliminare, intermedia o completa) dipende dal contesto.

\begin{regola}[Applicazione della Regola dell'Ottava senza numeri]
	\label{reg:applicazione_regola_ottava_basso}
	Quando il basso procede per gradi congiunti (scala ascendente o discendente), si applicano le regole seguenti:
	\begin{itemize}
		\item \textbf{I e VIII grado}: accordo di tonica in stato fondamentale (5).
		\item \textbf{II grado}: accordo sul VII in primo rivolto (6) oppure V4/3 --- a seconda del contesto e del livello.
		\item \textbf{III grado}: accordo di I in primo rivolto (6).
		\item \textbf{IV grado}: accordo di II in primo rivolto (6) oppure II6/5.
		\item \textbf{V grado}: accordo di V in stato fondamentale (5) oppure V7. In alternativa, I6/4 cadenzale o di passaggio.
		\item \textbf{VI grado}: accordo di IV in primo rivolto (6).
		\item \textbf{VII grado}: accordo di V in primo rivolto (6) oppure V6/5.
	\end{itemize}
	Per i dettagli e le varianti, si rimanda alla tabella completa in \ref{reg:regola_ottava_completa_maggiore}.
\end{regola}

\section{Riconoscimento delle progressioni nel basso senza numeri}
\label{sec:progressioni_basso_senza_numeri}

Le \textbf{progressioni} (Capitolo \ref{chap:progressioni_armoniche}) sono strumenti particolarmente utili nel basso senza numeri, perché permettono di riconoscere un pattern e applicare un sistema di armonizzazione coerente a tutto il passaggio.

\begin{regola}[Come riconoscere una progressione nel basso]
	\label{reg:riconoscimento_progressione}
	Una progressione si riconosce quando le note del basso formano un pattern che si ripete a distanza regolare:
	\begin{itemize}
		\item \textbf{Salti ripetuti di quarta o quinta}: indicano una progressione per quinte (ascendente o discendente).
		\item \textbf{Alternanza grado congiunto-salto}: può indicare una progressione per terze con rivolti alternati.
		\item \textbf{Discesa o salita cromatica}: indica una progressione cromatica (cfr. \ref{sec:progressioni_cromatiche}).
		\item \textbf{Scalare con ritmo regolare}: può indicare una progressione scalare secondo la Regola dell'Ottava.
	\end{itemize}
	Quando si riconosce una progressione, si applica lo stesso tipo di armonizzazione a tutte le ripetizioni del modello.
\end{regola}

\section{Alterazioni cromatiche: segnali di tonicizzazione e modulazione}
\label{sec:cromatismi_basso}

Nel basso senza numeri, ogni alterazione cromatica è un \textbf{segnale} che richiede un'interpretazione attenta. Le possibilità sono quattro, in ordine di complessità crescente:

\begin{enumerate}
	\item \textbf{Sensibile del modo minore}: l'alterazione è il VII grado alzato nella tonalità in corso. Non cambia nulla di strutturale; si usa semplicemente la dominante maggiore (V o V7) con la sensibile corretta.
	\item \textbf{Cromatismo ornamentale}: l'alterazione non ha funzione armonica, ma abbellisce la linea melodica del basso. Raro negli esercizi scolastici, ma da riconoscere.
	\item \textbf{Tonicizzazione} (cfr. \ref{chap:tonicizzazione}): l'alterazione crea una dominante secondaria che ``punta'' verso un grado non-tonicale come fosse una tonica locale. L'effetto è momentaneo e la tonalità principale si ristabilisce subito dopo.
	\item \textbf{Modulazione} (cfr. \ref{chap:modulazione}): l'alterazione segna l'ingresso in una nuova tonalità, confermata da una cadenza nella nuova chiave. Il cambiamento è strutturale e duraturo.
\end{enumerate}

\begin{regola}[Come interpretare un'alterazione nel basso senza numeri]
	\label{reg:interpretazione_alterazione}
	Di fronte a un'alterazione cromatica nel basso, chiedersi nell'ordine:
	\begin{enumerate}
		\item È il VII grado alzato del modo minore in corso? $\rightarrow$ sensibile, nessun problema.
		\item Questa nota è la \textit{terza} di una triade maggiore che risolve per quinta discendente sull'accordo successivo? $\rightarrow$ tonicizzazione o dominante secondaria.
		\item La tonalità sembra cambiare e viene confermata da una cadenza? $\rightarrow$ modulazione.
		\item Se nessuna delle precedenti: cromatismo ornamentale (usare la Regola dell'Ottava con la nota alterata come accordo di passaggio).
	\end{enumerate}
\end{regola}

\section{Conferma della tonalità: strumento di analisi e di composizione}
\label{sec:conferma_tonalita}

Sia nell'analisi che nella realizzazione del basso senza numeri, la \textbf{conferma della tonalità} è un concetto operativo fondamentale. Essa indica il processo con cui si stabilisce, senza possibilità di doppia lettura, la tonalità e il modo di un brano o di una sezione.

\begin{definizione}[Conferma della tonalità]
	\label{def:conferma_tonalita}
	La \textbf{conferma della tonalità} è la presenza di quegli accordi che, insieme, permettono di identificare inequivocabilmente la tonalità e il modo del brano. Il metodo più diretto è la \textbf{cadenza composta} (IV-V-I): i tre accordi costruiti su questi gradi contengono \textit{tutte} le note della scala, rendendo impossibile qualsiasi ambiguità tonale.
\end{definizione}

\begin{regola}[Quando e dove la tonalità viene confermata]
	\label{reg:quando_conferma_tonalita}
	\begin{itemize}
		\item \textbf{All'inizio del brano}: i primi accordi stabiliscono la tonalità. Anche solo I-V o I-IV sono sufficienti per orientare l'orecchio.
		\item \textbf{Dopo una modulazione}: la nuova tonalità deve essere confermata da una cadenza autentica (V-I) nella nuova chiave. Senza questa conferma, si parla di semplice tonicizzazione.
		\item \textbf{Alla fine del brano}: la cadenza finale (di solito composta II-V-I o IV-V-I) è la conferma definitiva della tonalità di impianto.
	\end{itemize}
\end{regola}

\section{Nota del basso che vale più tempi}
\label{sec:nota_piu_tempi}

Una delle situazioni più frequenti nel basso senza numeri è quella di una nota che vale più tempi, o addirittura più misure. In questo caso, lo studente deve decidere se assegnare un solo accordo oppure far cambiare l'armonia sopra la stessa nota del basso.

\begin{regola}[Nota del basso che vale più tempi]
	\label{reg:nota_piu_tempi}
	\begin{itemize}
		\item Una nota del basso che vale \textbf{più di un tempo} può appartenere a tonalità diverse o portare più di un accordo. La scelta dipende dal contesto fraseologico: se siamo in zona cadenzale, conviene armonizzarla con l'accordo cadenzale corrispondente.
		\item Una nota lunga in posizione di \textbf{tonica} porta di norma un unico accordo di tonica.
		\item Una nota lunga in posizione di \textbf{dominante} può portare I6/4 (cadenzale) risolto su V, oppure semplicemente V per l'intera durata.
		\item Se la nota lunga non è in posizione cadenzale, si può scegliere di armonizzarla con due accordi diversi, sfruttando ad esempio un cambio di posizione o un accordo vicino (per terza o per quinta).
	\end{itemize}
\end{regola}

\section{La nota del basso che salta}
\label{sec:nota_salta}

\begin{regola}[Nota del basso che salta: principio degli stati fondamentali]
	\label{reg:nota_salta_stati_fondamentali}
	Una nota di basso che è raggiunta per \textbf{salto} (non per grado congiunto) porta quasi sempre un accordo in \textbf{stato fondamentale}. I rivolti sono rari in corrispondenza dei salti, perché il salto segnala un cambio di pilastro armonico che richiede solidità strutturale.
	
	\textbf{Unica eccezione rilevante:} il I grado in primo rivolto (I6) può comparire dopo un salto di ottava ascendente quando il basso realizza un arpeggio della tonica.
\end{regola}

\begin{errore}[Accordo rivoltato dopo un salto nel basso]
	\label{err:rivoltato_dopo_salto}
	Evitare di assegnare accordi in secondo rivolto (6/4) o terzo rivolto (2) dopo un salto nel basso. Questi accordi instabili richiedono un contesto melodico che solo il grado congiunto può fornire.
\end{errore}

\section{Dalla deduzione alla realizzazione: il flusso di lavoro}
\label{sec:flusso_lavoro_basso_senza_numeri}

Il flusso di lavoro completo per il basso senza numeri si articola in queste fasi, che riprenderemo in modo dettagliato nel Capitolo \ref{chap:metodo_basso_senza_numeri}:

\begin{enumerate}
	\item \textbf{Analisi preliminare}: tonalità, modo, alterazioni cromatiche.
	\item \textbf{Segmentazione}: frasi, cadenze, struttura del periodo.
	\item \textbf{Classificazione di ogni nota}: grado della scala, posizione metrica, tipo di movimento.
	\item \textbf{Deduzione della numerica implicita}: applicare Regola dell'Ottava, progressioni, formule cadenzali.
	\item \textbf{Scrittura delle voci superiori}: applicare le regole di condotta delle voci conosciute.
	\item \textbf{Revisione}: checklist degli errori (\ref{reg:checklist_errori}).
\end{enumerate}

Il capitolo seguente (\ref{chap:metodo_basso_senza_numeri}) sviluppa ognuna di queste fasi in modo operativo, con esempi svolti commentati e strategie per la gestione dei casi più complessi.
